\section{Классификация протоколов многосторонних конфиденциальных вычислений}
\label{sec:protocol-classification}

\subsection{Классификация по модели противника}
\label{subsec:classification-threats}

Полубезопасная (semi-\allowbreak honest) модель. В этой модели предполагается, что нечестные участники следуют тексту протокола, но пытаются извлечь дополнительную информацию из получаемых сообщений.  Такой противник называется честный, но любопытный.  Формальное определение состоит в том, что коррумпированные стороны могут объединять свои виды, но не изменяют сообщения.  Протоколы в полу-честной модели предотвращают непреднамеренную утечку информации и часто служат основой для усиления гарантий; меньшее число криптографических проверок и обменов повышает эффективность.

Активная (malicious) модель. В активной (или злонамеренной) модели противник может отклоняться от протокола, подделывать сообщения, пытаться изменить вывод честных сторон.  Протоколы, обеспечивающие безопасность в этой модели, гарантируют, что при любых действиях противника честные стороны получают корректный результат, а приватность данных сохраняется.  Если при активном нарушителе невозможно вычислить результат (например, нарушитель пытается завершить вычисление преждевременно), честные участники либо обнаруживают факт нарушения и прерывают протокол, либо обеспечивается справедливость, когда без результата не остаётся ни одна честная сторона.  Защита от активного противника обычно использует доказательства корректности (zero-knowledge proofs), проверки совпадения входов и копии схем (cut-and-choose), увеличивая издержки.

Скрытная (covert) модель. Между полубезопасной и активной моделями выделяют промежуточный вариант, в котором предполагается, что противник готов нарушить протокол только если его действия останутся незамеченными.  Для таких скрытных adversaries протоколы обеспечивают обнаружение нарушения с заданной вероятностью (обычно 75---90 \%).  Модель covert применяется в коммерческих задачах, где репутационные или юридические последствия снижают вероятность атаки, и позволяет повысить эффективность относительно полной активной модели.

\subsection{Классификация по количеству участников}

Двусторонние протоколы (2PC). Двусторонний вариант МКВ предполагает, что только два участника \(P_A\) и \(P_B\) совместно вычисляют функцию \(f(x_A,x_B)\).  Для двух сторон применимы техники, не масштабируемые на множество участников.  Используется схема зашифрованных схем Яо, в которой функция описывается в виде булевой схемы, затем «garbler» отправляет «вычислителю» зашифрованную схему, после чего обе стороны получают результат, ничего более не узнавая.  Яо-протокол обеспечивает безопасность против полудобросовестного противника и характеризуется постоянным числом раундов общения, независимо от сложности функции.  Расширения этого подхода включают активную защиту с cut-and-choose и протокол BMR, который позволяет использовать garbled circuits в многостороннем случае.

Многосторонние протоколы (n-party). Когда вычисление выполняют более двух участников, роли garbler и evaluator исчезают; данные каждой проводки схемы представляются в виде долей, распределённых между сторонами.  В таких протоколах функция изображается как арифметическая схема над конечным полем, а операции сложения и умножения выполняются над секретными долями.  Типичный пример --- протокол GMW, разработанный Голдрайхом, Микали и Вигдерсоном; он использует аддитивное разделение секрета и масштабируется на произвольное число сторон.  GMW допускает активных нарушителей, если большинство участников честное, и вводит доказательства с нулевым разглашением для защиты от обманов.  Другой пример --- протокол BGW, который строится на схеме Шамира и использует арифметические схемы.  Он эффективен для функций, описываемых операциями сложения и умножения, и позволяет устанавливать пороговое значение количества нечестных участников.

Допущение о честном большинстве. В протоколах на основе разделения секрета важную роль играет порог \(t\) числа нечестных сторон.  Для схемы Шамира секрета величина порога, при котором обеспечивается информационно-\allowbreak теоретическая конфиденциальность, зависит от модели противника: при пассивном нарушителе условие \(t< n/2\) достаточно для восстановления секрета, а при активном нарушителе требуется \(t< n/3\).  Аддитивное разделение секрета может выдерживать до \(n-1\) нарушителей в обеих моделях, но требует дополнительных механизмов для проверки корректности (например, MAC-аутентификации в протоколе SPDZ).  Таким образом, классификация по количеству участников тесно связана с предположением об относительном числе честных сторон, что влияет на выбор схемы и гарантии безопасности.

\subsection{Классификация по используемому методу}
\label{subsec:classification-primitives}

В зависимости от используемых криптографических примитивов протоколы МКВ подразделяют на несколько групп. Механизмы и сравнительная характеристика примитивов приведены в разделе~\ref{sec:cryptographic-foundations}. Критерии согласуются с обзором П. Морейры, сопоставляющим протоколы по предположениям безопасности, эффективности, удобству использования и функциональным возможностям \cite{10}.

Протоколы на основе разделения секрета используют секретные доли для представления данных. К этой группе относятся BGW, GMW и SPDZ. Гарантии безопасности каждого протокола зависят от принятой модели противника.

К протоколам на основе зашифрованных схем относятся протокол Яо и его многосторонние расширения, например BMR. Отдельную группу составляют протоколы на основе гомоморфного шифрования, в которых данные обрабатываются посредством операций над шифротекстами.

Гибридные протоколы сочетают несколько криптографических примитивов на разных этапах вычисления. Классификация по основному методу не исключает вспомогательных механизмов, в частности ослеплённой передачи (см. раздел~\ref{sec:cryptographic-foundations}).

Специализированные протоколы для приватного пересечения множеств, аукционов и других задач выделяют по назначению, а не по используемому криптографическому методу. Они могут относиться к любой из перечисленных групп. Обзор PSI-протоколов Д. Моралеса, И. Агудо и Х. Лопеса выделяет гомоморфное шифрование, ослеплённую передачу и забывающие псевдослучайные функции и сопоставляет решения по сложности и практической применимости \cite{09}. Классификация по назначению дополняет, но не заменяет критерии примитивов, числа участников и модели противника.

\subsection{Вывод}

Критерии выбора протокола заданы классификацией (см. подразделы~\ref{subsec:classification-threats}--\ref{subsec:classification-primitives}). Двусторонние схемы (garbled circuits, FHE) обеспечивают низкую латентность при сравнении и обработке конфиденциальных данных двух сторон; многосторонние GMW, BGW и SPDZ применяются для распределённых вычислений при честном большинстве. Зависимость коммуникационных и вычислительных затрат от примитивов отражена в таблице~\ref{tab:mpc-primitives}. В следующем разделе рассмотрены применения МКВ в аналитике данных, МО, финансах и медицине.